\documentclass[12pt]{article}

\usepackage[T1]{fontenc}
\usepackage[utf8]{inputenc}
\usepackage[brazil]{babel}
\usepackage{lmodern}
\usepackage{microtype}
\usepackage[a4paper,margin=2.5cm]{geometry}
\usepackage{setspace}
\usepackage{indentfirst}

\setstretch{1.15}

\title{Política Econômica na Economia Marginalista e Neoclássica Inicial, 1871–1920s \thanks{Texto Traduzido livremente por Luiz Mario Andrade com o auxílio da ferramenta Chat GPT 5.4}}
\author{Michel S. Zouboulakis}
\date{2023}

\begin{document}

\maketitle

\section*{Introdução}

A economia depois de 1871 realizou uma guinada epistemológica significativa devido à chamada “Revolução Marginalista”, que é geralmente associada à descoberta simultânea do conceito de utilidade marginal e à aplicação do cálculo diferencial em vários campos da teoria econômica. Ambas as inovações foram possíveis graças ao deslocamento do escopo da análise econômica, das condições sociais de produção, distribuição e acumulação de riqueza, para os fatores subjetivos de troca e produção, ou seja, da Economia Política para a Ciência Econômica (Economics). À medida que a disciplina se tornou mais formal e mais individualista, ela se tornou mais estreita em seu escopo. Como observou Heilbroner (1953, 173), “tornou-se a província especial de professores, cujas investigações lançavam feixes pontuais em vez dos faróis de busca ampla dos economistas anteriores”.

O período 1873–1896 é marcado pela Longa Depressão, durante a qual a economia mundial enfrentou preços em queda constante, taxas de juros e de lucro em declínio, gerando falências empresariais massivas e aumento do desemprego.\footnote{Apesar da visão estabelecida, não houve uma depressão real após 1873, mas apenas uma desaceleração. Assim, para o principal fabricante mundial, a taxa de crescimento anual da produção industrial britânica caiu de 2,7 em 1853–1873 para 2,1 em 1873–1899 e a taxa de crescimento do PIB total de 1,95 para 1,85, respectivamente (Lewis 1978, 112–113).} O crescimento econômico sustentado retornou lentamente após a reorganização substancial do trabalho e do capital, o surgimento do protecionismo e o empurrão da expansão imperialista aos seus limites. Como a virada para a teorização formal afetou a visão do economista sobre questões candentes de política econômica? Este capítulo considera concisamente as opiniões sobre política econômica de algumas mentes econômicas excepcionais, como W.S. Jevons, L. Walras, A. Marshall, F.Y. Edgeworth, J.B. Clark, V. Pareto, K. Wicksell, I. Fisher e A.C. Pigou (para citar os mais influentes), cobrindo quase 60 anos, de 1871 a 1929, em três grandes questões de política: comércio internacional, política monetária e finanças públicas. Será mostrado que os economistas tradicionais daquele período não estavam menos envolvidos com considerações sociais e políticas do que seus predecessores clássicos e não eram menos heterogêneos quanto às suas prescrições de política.

\section*{Política de Livre Comércio da Grande Depressão aos “Anos Vinte Rugientes”}

A maior parte do período em estudo foi chamada por Hobsbawm (1987) de “A Era dos Impérios, 1875–1914” por uma razão: ela corresponde à era da maior expansão colonial da Europa no exterior e à escalada do nacionalismo em todo o mundo.\footnote{Em 1913, os poderes coloniais europeus dominavam mais de 550 milhões de pessoas, o que é 2,7 vezes a população europeia (Bairoch 1997, vol. II, 557), embora eles se expandissem após 1918 para os ex-territórios otomanos do Oriente Médio.} O surgimento do protecionismo foi apenas uma consequência do crescente antagonismo entre as nações. Como observou Hobsbawm, em 1914 o mundo da sociedade burguesa liberal havia desaparecido irreversivelmente (Hobsbawm 1987). É, portanto, nesse clima particularmente antiliberal que os economistas marginalistas e neoclássicos iniciais desenvolveram suas recomendações de política sobre o comércio internacional.

Lendo os manuais de Economia Política das décadas de 1870 e início de 1880 na Grã-Bretanha (ou seja, Mill, Cairnes, Fawcett e Sidgwick), compreende-se que a academia ainda endossava totalmente a teoria clássica do comércio internacional. Em poucas palavras, Mill aprimorou a teoria ricardiana da vantagem comparativa explicando que o intercâmbio internacional depende dos termos de troca entre as razões de custo comparativo dos dois países. A demanda recíproca, ou “a quantidade de produção doméstica que deve ser dada ao país estrangeiro em troca dela”, determina o preço da mercadoria importada (Mill 1848, 584). Na Grã-Bretanha e na França, economistas e políticos, como Mill e Gladstone, celebraram o sucesso do Tratado Cobden-Chevalier de 1860, que encerrou as tarifas sobre os principais itens de comércio entre as duas superpotências econômicas. A Grã-Bretanha assinou também tratados comerciais com a Bélgica (1862), o Zollverein (1862), a Itália (1863), a Suíça (1864), a Espanha e a Holanda (1865). Jevons (1835–1882) em seu \textit{The Coal Question} (1865) afirmou que manter preços baixos de carvão importado era crucial para uma nação industrial em crescimento manter sua competitividade e prosperidade (Mosselmans 2007, 14). Mas, à medida que a industrialização avançava e a posição econômica da Alemanha crescia após a Unificação de 1871 e a economia dos EUA revelava sua força após a Guerra Civil, a agenda política protecionista tornou-se popular novamente em meados da década de 1880, nos piores dias da Grande Depressão.\footnote{Karl Knies, Gustav Schmoller na Alemanha, e Henry Carey, R.T. Ely nos EUA, difundiram ideias protecionistas.} Pode-se ver a evolução no espírito da época através das seis edições consecutivas de \textit{Free Trade and Protection} de Henry Fawcett, com o objetivo de dar “consideração completa e cuidadosa aos argumentos avançados pelos protecionistas” (1872, 12), especialmente aqueles relacionados ao movimento de “comércio justo”. Recompensas à exportação, subsídios protetores e várias restrições às importações são vantajosos para um país apenas se este país desejar “impedir a prosperidade dos países vizinhos”, iniciando assim uma guerra econômica perpétua (1872, 73). Como conclui Fawcett (192), não se deve considerar apenas os interesses dos fabricantes, ignorando assim “o interesse do povo, como consumidores”. Argumentos semelhantes foram avançados na América por Henry George em seu \textit{Protection or Free Trade}, publicado em 1886: as tarifas apenas mantêm os preços altos para os consumidores, falham em produzir qualquer aumento nos salários e, no final do dia, aumentam o poder das empresas monopolistas. No mesmo ano, John Bates Clark (1847–1938), em sua \textit{Philosophy of Wealth}, fez a defesa analítica mais forte do livre mercado e dos limites da intervenção governamental. Clark (1886, 1901) defendeu a ampla concorrência contra práticas de preços predatórios, tanto nacional quanto internacionalmente, e contribuiu significativamente para a propagação da legislação antitruste nos EUA (cf. Pressman 2014, 105).

Alfred Marshall (1842–1924) é o primeiro grande economista convencional que realmente deu algum crédito ao novo humor protecionista. Em seus \textit{Principles of Economics} (1890, 633), ele reconhece que o livre comércio era mais vantajoso para a Grã-Bretanha do que para a Alemanha ou os EUA, juntamente com o antigo argumento milliano sobre a necessidade de proteção temporária de um país ou setor industrial jovem. Mas, embora compreenda o significado patriótico das medidas protecionistas sugeridas pela Escola Histórica Alemã, ele não aprova suas visões. A posição de Marshall evoluirá durante a controvérsia da reforma tarifária em 1903, quando ele “permitiu que seu nome fosse incluído no manifesto de quatorze professores publicado no \textit{The Times}” em favor de tarifas protetoras (Backhouse, em Raffaelli et al. 2006, 478). Após o fim da Primeira Guerra Mundial, seu \textit{Industry and Trade} soa mais compreensivo ao “nacionalismo econômico”, reconhecendo a eficiência das tarifas protetoras para a indústria nascente em países em desenvolvimento: os economistas ingleses “negligenciaram o fato de que muitos daqueles efeitos indiretos da Proteção que se agravaram então, e se agravariam agora, seus males diretos na Inglaterra, trabalharam na direção oposta na América” (Marshall 1919, 476). Após um relato histórico minucioso da história econômica da França, Alemanha, EUA e outros países como o Japão, Marshall reconhece que a “liderança na indústria e no comércio foi obtida no passado por cidades ou por nações que lançaram energia no uso de suas próprias 'mercadorias nativas'” (1919, 76).

Em suma, em ambos os lados do Atlântico, a melodia popular era protecionista à medida que nos aproximávamos da Primeira Guerra Mundial e a ambivalência de Marshall reflete a pressão da opinião pública. No entanto, ele nunca abandonou o cânone do livre comércio, acreditando que a proteção estatal se torna prejudicial ao comércio, pois muitas vezes se torna um instrumento de tratamento preferencial e corrupção política (Marshall 1923). Apesar do auge do imperialismo e do nacionalismo, os economistas neoclássicos permaneceram defensores do livre comércio em princípio e na teoria.\footnote{Além do dogma liberal, F. Y. Edgeworth (1894) usará curvas de oferta e curvas de indiferença comunitárias para explorar a questão da tarifação ótima.} Como Heilbroner (1953, 193) brilhantemente concluiu, “o oficialato da economia ficou de um lado, observando o processo de crescimento imperial com equanimidade e limitando suas observações ao efeito que as novas possessões poderiam ter sobre o curso do comércio”.

\section*{Política Monetária e as Origens do Monetarismo}

A segunda questão de política mais importante dizia respeito ao dinheiro e ao padrão-ouro. No Reino Unido, após o \textit{Bank Charter Act} de 1844, as ideias da Escola de Moeda (Currency School) prevaleceram apesar da monumental \textit{History of Prices} de Thomas Tooke, demonstrando a independência do movimento dos preços em relação à quantidade circulante de dinheiro.\footnote{Thomas Tooke publicou seus seis volumes de \textit{History of Prices} de 1838 a 1857 — os dois últimos em coautoria com William Newmarch — para demonstrar empiricamente que os preços variavam ao longo do tempo principalmente devido a boas colheitas, variações cambiais estrangeiras, obstruções ao livre comércio, preços altos de matérias-primas, inovações tecnológicas, etc. Ver Schumpeter (1954, 708–713).} A Economia Política Clássica era dedicada à chamada teoria quantitativa da moeda e, na verdade, concedeu essa ideia bicentenária (desde Petty e Locke) — de que os preços das mercadorias dependem da quantidade do meio circulante — à geração seguinte. Juntamente com Mill e Jevons, Marshall manteve essa posição desde seu primeiro “Essay on Money” em 1871 até seu último livro de 1923 e afirmou que:

\begin{quote}
se tudo o mais permanecer igual, existe esta relação direta entre o volume da moeda e o nível de preços, de modo que, se um for aumentado em dez por cento, o outro também será aumentado em dez por cento. (Marshall 1923, 45; Cf. Mill 1848, 495)
\end{quote}

Juntamente com a equiproporcionalidade, Marshall expõe todos os elementos da teoria quantitativa clássica, como a neutralidade da moeda, a exogeneidade, a dicotomia preço absoluto/preço relativo e a causalidade indireta da moeda para o preço (Humphrey, em Raffaelli et al. 2006, 421–433). Marshall escreveu: “O dinheiro não é desejado por si só” (1923, 38), mas, ao mesmo tempo, ele também reconheceu que o dinheiro em espécie pode não ser inteiramente gasto, mas mantido para “gastos futuros” (1890, 189), criando assim um elo entre o crescimento do dinheiro e sua circulação.

Em questões de política monetária, nota-se esta impressionante continuidade teórica de Mill para Jevons e Marshall: a regulação da oferta monetária era de inteira responsabilidade do Estado. A norma era um sistema monetário baseado em mercadorias com conversibilidade total em um metal precioso (ouro e/ou prata). Não obstante o acordo geral para um sistema bancário livre competitivo, Walter Bagehot (1826–1877) em seu ilustre \textit{Lombard Street} (1873) descreveu a obrigação do banco central de manter as reservas adequadas dos bancos comerciais, mas também de manter a solvência da nação e preservar a viabilidade de todo o sistema financeiro. Bagehot queria que o Banco da Inglaterra (BoE) “reconhecesse sua posição especial na base da pirâmide de crédito da economia” (Laidler 1991, 42) e afirmou que o BoE tinha o dever de emprestar livremente aos banqueiros comerciais a taxas elevadas, quando nenhum outro credor o faria, muito antes de Keynes.\footnote{Assim, Keynes escreveu sobre ele como sendo “o livro em toda a biblioteca de literatura econômica que todo estudante de economia, por mais humilde que seja, terá lido, embora não leia mais nada” (EJ, 1915, 371). \textit{Lombard Street} foi um enorme sucesso (21 edições de 1873 a 1933), tornando Bagehot a maior autoridade em bancos centrais. Taxas elevadas são sugeridas para evitar que os banqueiros tomem dinheiro emprestado sem precisar dele.} Ainda assim, Bagehot distinguiu-se da ortodoxia monetária da Escola de Moeda questionando a causalidade entre a oferta de dinheiro e as flutuações de preços, acreditando que a quantidade de dinheiro era uma variável passiva que seguia as flutuações de preços em vez de causá-las.

E, no entanto, a ortodoxia monetária britânica não era monetarista! O monetarismo veio mais tarde com Irving Fisher (1867–1947) e a exposição formal definitiva da teoria quantitativa da moeda em \textit{The purchasing power of money} (1911). Diretamente inspirado pela Lei de Boyle para Gases Perfeitos (cf. Zouboulakis 2003), e seguindo inúmeros estudos históricos e empíricos sobre a história monetária dos EUA, Fisher reescreve a equação de Simon Newcomb desejando explicar o poder de compra da moeda.\footnote{Fisher refere-se às obras de Simon Newcomb (1885), Alexander Del Mar (1885), Charles Dunbar (1901), Charles Laughlin (1903), Wesley Mitchell (1903) e principalmente Edwin Kemmerer (1909). Ver Tobin (1985) e Laidler (1991).} Indo além da literatura existente na época, Fisher deseja estabelecer uma relação causal entre o nível de preços (P) e a quantidade de dinheiro em circulação (M). Ao fazê-lo, ele considera que a velocidade de circulação da moeda (V) é determinada por instituições e hábitos e muda imperceptivelmente ao longo do tempo e que, no curto “período de transição”, a quantidade de bens (T) comprados por M “dentro de um ano” é dada. Sua equação de troca $MV = PT$ não pretendia expressar uma condição de equilíbrio (como Walras), mas explicar a causalidade unidirecional do aumento exógeno da oferta de moeda para os preços nominais.\footnote{Fisher expressou-a pela primeira vez como $MV = \sum pQ$, ou seja, a soma dos preços multiplicada pelas quantidades compradas (1911, 26) e depois como $MV + M'V' = \sum pQ$, para incluir a quantidade de depósitos em conta corrente e sua velocidade separada (1911, 149).} Como ele escreveu, “o nível de preços é normalmente o único elemento absolutamente passivo na equação de troca” (Fisher 1911, 172). Embora Fisher tenha oferecido uma verificação estatística detalhada dos ajustes transicionais entre os diferentes elementos da equação (1911, cap. XI–XII. Cf. Laidler 1991, 81–82), seu principal objetivo era proclamar seu programa monetarista em alto e bom som, considerando V e T como estáveis e M como exógeno e independente. Schumpeter (1954, 1103) conjecturou acertadamente que Fisher sacrificou a adequação científica pela utilidade prática: a teoria quantitativa foi “concebida como um andaime para o trabalho estatístico que, por sua vez, serviria a uma peça de engenharia social”. Seu “plano do dólar compensado” visava restringir a emissão de novo dinheiro dos bancos comerciais a 100\% de suas reservas de ouro. A teoria de Fisher era meramente o respaldo teórico de sua cruzada para estabilizar o valor do dólar (Tobin 1985, 35). Nesse sentido, divergiu não apenas da antiga tradição quantitativista britânica, mas também da nova abordagem alternativa, a teoria do Saldo de Caixa de Cambridge (ver Blaug 1985, cap. 15).

A teoria quantitativa central para as políticas monetaristas de Fisher, e mais tarde de Friedman, foca na oferta de moeda, considerando-a como o principal motor do crescimento econômico e, no longo prazo, como a principal causa da inflação se expandida indevidamente. Portanto, a oferta de moeda deveria ser a principal preocupação dos tomadores de decisão de política econômica. Pelo contrário, Mill, Bagehot e Marshall consideram também a demanda por moeda e os motivos para reter moeda, abrindo caminho para o conceito de “preferência pela liquidez” de Keynes (Mill 1848, 496; Bagehot 1873; Marshall 1890, 189. Cf. Raffaelli 2003, 134). Os economistas de Cambridge depois de Marshall, especialmente Pigou e Keynes, expressaram a demanda por saldos de caixa dependendo da proporção da renda que as pessoas desejam manter na forma de dinheiro (k), o nível de preços (p) e o nível de renda real (Y): $M = k \cdot p \cdot Y$. Nesse sentido, as pessoas podem não desejar usar o dinheiro como espécie imediatamente e armazená-lo para uso posterior, reduzindo assim sua quantidade em circulação. O fator crucial é a taxa de juros que remunera a abstinência do consumo imediato.

Knut Wicksell (1851–1926) foi o economista que desenvolveu completamente a teoria dos juros e assegurou o vínculo entre as políticas monetária, bancária e de estabilização. Ele distingue o “juro de mercado”, a remuneração pela poupança, do “juro natural”, o retorno (ou eficiência marginal) do novo capital investido (Wicksell 1898).\footnote{A distinção entre 'taxa de mercado' e 'taxa natural' já estava presente em Mill (1848, 639–640), dentro de um quadro teórico clássico.} A margem entre os dois juros determina a decisão de poupar ou investir. Como os bancos têm a possibilidade de estender o crédito além de seus depósitos em dinheiro, mantendo a lacuna entre os dois juros, eles podem alimentar a demanda do consumidor por bens acima de sua taxa natural e elevar os preços. Wicksell desafia, assim, a teoria quantitativa clássica da moeda muito antes da \textit{Teoria Geral} de Keynes, demonstrando teoricamente que a demanda agregada pode ser artificialmente sustentada por juros bancários inferiores à taxa normal, criando um efeito cumulativo sobre os preços (Wicksell 1898, 95, 1907). Do ponto de vista da oferta agregada, o juro bancário pode causar depressão econômica, ao manter a taxa de juros de mercado acima da taxa natural. Nesse caso, o investimento torna-se muito caro e a produção cai. Wicksell oferece, pela mesma razão, uma explicação substancial para as flutuações econômicas, a divergência entre as taxas de juros natural e de mercado, em oposição à explicação monetarista de Fisher sobre os ciclos de negócios causados pela “ilusão monetária”, ou a divergência entre o juro nominal e o real. A sugestão de Wicksell foi desenvolvida por Ludwig von Mises em seu primeiro trabalho \textit{The Theory of Money and Credit} (1912), antes de se tornar a pedra angular da Escola de Economia de Estocolmo na década de 1930.

\section*{Finanças Públicas e o Papel do Estado até 1929}

Em relação à política econômica interna e ao papel do Estado, os primeiros economistas neoclássicos acreditavam que “a função do Governo é governar o menos possível, mas não fazer o menos possível” (Pigou 1925). Inspirados pelos últimos economistas políticos clássicos, como Mill, Cairnes e Sidgwick, os Marginalistas ampliaram as funções tradicionais do Estado (segurança interna e externa, justiça, educação obrigatória, infraestruturas) usando ferramentas diferentes e a partir de um ponto de partida metodológico distinto. Depois de Marshall, a maioria dos economistas neoclássicos opôs-se à distinção entre a ciência e a arte da economia, que era a marca registrada da escola clássica depois de Ricardo e mesmo em Walras (1874, 39). Para Marshall, não havia ciência pura e toda a economia era arte, desafiando a opinião de seu pai espiritual, Henry Sidgwick.\footnote{Em seu \textit{Principles of Political Economy} (1883, 13, 395–396) Henry Sidgwick ainda subscreve a ideia clássica da distinção entre a ciência e a arte da economia política estabelecida por Senior e Mill, meio século antes. Os dois ramos diferem em seu escopo e no método de raciocínio, sendo a ciência positiva e hipotética e a arte sendo a prática voltada para obter “algum resultado desejável não para um indivíduo, mas para uma comunidade política”. Foi Ronald Coase quem descreveu Sidgwick como a “mãe e pai espiritual” de Marshall.} Consequentemente, o sucessor de Marshall em Cambridge, Arthur Cecil Pigou (1877–1959), afirmou que estava interessado na ciência que dá frutos (fruit-bearing), não na ciência que traz luz (light-bearing) (1920, 3). A economia pura separada da política, como Walras sonhava, não tinha significado, pois a economia é concebida como “a base de uma arte” que visa aumentar o bem-estar social. Para esse propósito, o Estado desempenha um papel significativo na realocação de recursos por meio de impostos e subsídios, a fim de curar as falhas do mercado que impedem o impulso do “dividendo nacional” e diminuem o bem-estar de uma nação. Entre sua análise sistemática da intervenção estatal, a contribuição mais conhecida de Pigou refere-se aos mecanismos de redução do custo social por meio da internalização de externalidades.

O conceito de externalidade foi discutido pela primeira vez por Sidgwick como resultado da observação de que os atores individuais falham em estar cientes do impacto social total de suas atividades (Sidgwick 1883, 412–413. Cf. Medema 2009, 43). Em seguida, Marshall referiu-se a externalidades positivas (sem usar o termo) como um corolário da ideia de economias externas produzidas pela concentração de indústrias especializadas em localidades particulares (1890, Livro IV, cap. x–xi; 1919, 287). Marshall compreendeu perfeitamente a ideia de impulsionar externalidades positivas e sugeriu uma “taxa de ar fresco”, um imposto moderado especial sobre terras urbanas com “valor de local especial” a ser gasto pelas autoridades locais “em abrir pequenos pontos verdes em meio a distritos industriais densos” (Marshall 1890, 662; Cf. Caldari e Masini 2011, 723). No entanto, foi Pigou quem formulou a externalidade como um problema de ineficiência de mercado e sugeriu como remédio a realocação de recursos pelo Estado ou alguma autoridade pública, e “colocou as ideias de Sidgwick no quadro teórico marshalliano” (Medema, em Raffaelli et al. 2006, 640). Pigou (1920, 174) concebeu as externalidades produtivas de forma muito semelhante aos livros didáticos atuais, como uma divergência entre os valores do “produto líquido privado marginal e do produto líquido social marginal”. A fim de evitar que um terceiro seja cobrado com “disserviços” extras, diminuindo assim o “produto líquido social”, o governo pode remediar externalidades negativas via impostos, a menos que um acordo seja concluído entre as partes envolvidas. A contribuição mais famosa de Pigou para a economia, o “imposto pigouviano”, não era uma panaceia, mas apenas uma cura eficaz, quando a compensação direta ou remédios legais eram ineficazes para melhorar o bem-estar econômico (cf. Collard, em Raffaelli et al. 2006, 594; Pressman 2014, 136). Resumidamente, Pigou argumentou que, como a empresa industrial não tem incentivo para internalizar o custo social marginal que produz, o Estado deve impor um imposto sobre essa empresa que corrigiria o resultado ineficiente do mercado. Do ponto de vista do mercado, na presença de externalidades, os preços não expressam o custo total de produção e a empresa produz em excesso, desperdiçando recursos raros.

Vilfredo Pareto (1848–1923) também visava oferecer as ferramentas para avaliar o bem-estar econômico de forma objetiva. Em contraste com o bem-estarismo de Pigou, Pareto (1909) introduziu medidas ordinais para estimar a alocação ótima de recursos, contornando assim o impasse da medição da utilidade e das comparações interpessoais de utilidade (cf. Drakopoulos 1989, 40). Partindo de suas posições marginalistas iniciais e do legado de Walras, Pareto concentrou-se não no mapa de gostos individuais, como Edgeworth fez, mas sim nas ações observáveis para adquirir bens. Seu célebre critério de otimalidade ajuda a decidir se uma medida de política é boa ou má: boa, se beneficia pelo menos alguém sem prejudicar ninguém; má, se prejudica pelo menos um sem beneficiar ninguém. No entanto, a otimalidade de Pareto também repousa em suposições que são definitivamente carregadas de valor, como a ideia de que o bem-estar social é definido como a soma do bem-estar dos indivíduos conforme eles o percebem (Blaug 1985, 591–592).

Backhouse et al. (2021, 1) argumentaram recentemente que “quando os economistas abordaram problemas práticos, eles adotaram uma gama muito mais ampla de julgamentos éticos além do bem-estarismo”, significando uma abordagem às decisões sociais que reivindica o bem-estar individual como o principal critério de bem-estar social. Inesperadamente, isso se aplica também ao caso dos primeiros neoclássicos. Marshall declarou que “a riqueza existe apenas para o benefício da humanidade. Ela não pode ser medida adequadamente em jardas, nem mesmo como equivalente a tantas onças de ouro; sua verdadeira medida reside apenas na contribuição que faz para o bem-estar humano” (Pigou 1925, 366). Analogamente, Pigou (1920, 11) afirmou que o bem-estar econômico é “aquela parte do bem-estar social que pode ser trazida direta ou indiretamente em relação com a régua de medição do dinheiro”. Apesar dessa afirmação, o próprio Pigou reconheceu que os gostos individuais dependem também das condições sociais, tornando assim suas escolhas baseadas em critérios imprecisos de comparação. Portanto, os dividendos sociais poderiam ser usados apenas como um guia aproximado para avaliar as políticas de redistribuição dos ricos para os pobres. Pareto (1897) reconheceu francamente os limites da análise econômica e não apenas instou os economistas a considerar “outros modelos de homem mais próximos da realidade” quando se trata de aplicar suas teorias à prática, mas também reduziu o campo da análise econômica apenas a “ações lógicas”. A maior parte do comportamento social — o domínio das ações não-lógicas — tornou-se o objeto de estudo da Sociologia, e os economistas tiveram que levar em conta suas conclusões para fazer sugestões de políticas (Pareto 1909, 41, 1916).\footnote{Ações não-lógicas não eram necessariamente irracionais, mas apenas ações que seguiam um processo de deliberação não-lógico (Cf. Steiner 1995; Zouboulakis 2014, 40).}

Em muitos campos específicos da política econômica, os Marginalistas e os primeiros neoclássicos esposaram visões divergentes, pois compartilhavam diferentes ideologias políticas. O caso de Marshall é muito instrutivo. Groenewegen (1995, cap. 16) descreve sua “tendência ao socialismo” como sendo sempre do tipo moderado, apoiando a regulação e intervenção estatal para alcançar a “igualdade de oportunidades”, mas hostil à extensão indevida da propriedade pública e aos métodos burocráticos. O papel de Marshall como membro da Comissão do Trabalho de 1891 a 1895 revela um economista particularmente envolvido com medidas de política social visando o alívio da pobreza e a saúde (Jensen 1987). Sua atitude evoluiu nas edições subsequentes de seus \textit{Principles}, como mostram suas propostas para promover “o bem-estar da classe trabalhadora mais pobre” por meio da regulação estatal de “que nenhuma fileira de edifícios altos seja erguida sem espaço livre adequado na frente e atrás” (Marshall 1890, 597–598).\footnote{Marshall endossou as visões pró-unionistas de Mill contra a ideia de Edgeworth de que os sindicatos criam indeterminação na negociação salarial. Ver Creedy (1990). Da mesma forma, Clark sugeriu que os sindicatos estão dificultando a concorrência (Pressman 2014, 104).}

Igualmente notável foi o caso de Leon Walras (1834–1910). Seu excelente modelo de Equilíbrio Geral descreveu a economia como um único mercado gigante compreendendo todos os bens, fatores de produção e serviços em uma operação comercial unida. Este modelo abstrato tinha a qualidade — entre outras coisas — de divulgar a ideia de que diferentes setores de uma economia estão inter-relacionados. Mas quando se trata de falar sobre a aplicação deste modelo a questões de política, Walras abandona o regime hipotético de concorrência perfeita e torna-se bastante crítico das políticas de “laissez-faire”: “Obviamente! Deveria ser provado que a livre concorrência de fato maximiza a utilidade” (1898, 418). A principal razão para isso reside em sua tripla divisão entre “économie pure”, “économie sociale” e “économie appliquée”. As duas últimas subdisciplinas não dependem das leis da natureza como a economia pura, mas da “vontade humana que é uma força livre e perspicaz” (Walras 1874, 39).\footnote{Walras reconheceu que J.S. Mill (1848, 200) e sua distinção entre leis de produção e leis de distribuição foram a fonte de sua inspiração.} Como criações humanas, os fenômenos econômicos aplicados e sociais extraem seus princípios não exclusivamente da ciência, mas também da Moral. Um princípio de política, como o “laissez-faire”, deve ser avaliado de acordo com sua eficiência prática e impacto no bem-estar social. Igualmente, a intervenção estatal tem a responsabilidade de fornecer bens e serviços de interesse público e usar sua autoridade para reduzir a desigualdade social e promover a liberdade individual (Walras 1896, 162. Cf. Dockès 1996). Pelo contrário, a ideia de igualdade social era totalmente estranha ao seu sucessor em Lausanne (Pareto 1900. Cf. Legris 2000). Pareto sugeriu uma reformulação do processo de seleção natural para explicar a desigualdade social como uma fase da ordem social. Além disso, a presença da elite no topo da hierarquia social é uma indicação de uma seleção social eficaz (Pareto 1909, §97–115). Nessa base, qualquer intervenção estatal para proteger um grupo social particular perturbaria o equilíbrio social e produziria o efeito oposto, a longo prazo, favorecendo os membros mais fortes da sociedade (Pareto 1909, §2216–2217. Cf. Steiner 1995).

Discordâncias também se manifestam em relação ao eterno problema da tributação ótima: como cobrar um nível de impostos de modo a maximizar as receitas públicas com o menor dano tanto à produção quanto ao consumo, seguindo as quatro máximas de Adam Smith.\footnote{Em seu Capítulo 2 do último Livro de \textit{The Wealth of Nations}, Adam Smith formulou quatro máximas em relação aos impostos: equidade, certeza, conveniência de pagamento e economia de cobrança.} Em certo sentido, cada imposto contém alguma injustiça, pois deduz parte da renda pessoal obtida por meio do trabalho, investimento ou empreendimento. Assim, Mill (1848, 808) aceitou o papel redistributivo da tributação direta até certo limite: “não para aliviar o pródigo às custas do prudente”. Marshall concordou com isso como ponto de partida e apoiou o imposto de renda progressivo, exatamente porque ele diminui a desigualdade social excessiva, embora tenha reconhecido seus efeitos negativos sobre o excedente do consumidor (Marshall 1890, 661. Cf. Reisman em Raffaelli et al. 2006, 497). Walras e Pareto, pelo contrário, eram absolutamente contra qualquer imposto de renda progressivo por razões diferentes. Walras acreditava que o imposto de renda desafia o princípio da liberdade natural e aceita apenas impostos sobre a propriedade (Walras 1896, 413. cf. Potier 2019). Pareto, diretamente inspirado por Frederic Bastiat, considerava que os impostos consistem em uma forma de “espoliação” de riquezas pelo Estado e, portanto, deveriam ser reduzidos a um mínimo (Pareto 1909, §26–27).

A ideia mais radical sobre tributação veio de um economista político e ativista autodidata americano, Henry George (1839–1897). Acreditando ter descoberto a principal causa da pobreza, que era também o principal impedimento ao progresso social, George sugeriu um imposto único massivo sobre a fortuna da propriedade da terra que substituiria outros impostos e forneceria serviços sociais e também financiaria serviços públicos. Embora \textit{Progress and Poverty} carecesse das qualidades analíticas dos livros didáticos de Economia Política daquela época, foi vendido em milhares de exemplares e tornou-se o \textit{best-seller} por quase meio século (Heilbroner 1953, 188–190). Tanto as ideias de um imposto único quanto o impacto negativo das rendas não eram novos. Desde 1760, Quesnay e depois Turgot defenderam a necessidade do \textit{impôt unique} sobre a propriedade da terra, e Ricardo, desde seu ensaio de 1815 e definitivamente em seus \textit{Principles} de 1817, demonstrou o impacto dramático das rendas na acumulação capitalista e o futuro sombrio da economia estacionária. No entanto, George iniciou um movimento de reforma social — o Georgismo — que influenciou muitas pessoas, desde os Fabianos até os Institucionalistas Originais, como J.R. Commons.

Para contextualizar o debate acima, deve-se mencionar que, até a Primeira Guerra Mundial, as receitas públicas nos países ocidentais desenvolvidos, provenientes de qualquer fonte de tributação, representavam menos de 10\% do seu PIB (Piketty 2013, gráfico 13.1). Uma grande parte dessas receitas continuou a vir de tarifas sobre bens importados que cresceram rapidamente na Europa após 1880. Com exceção da Grã-Bretanha, Holanda e Dinamarca, todas as nações europeias aumentaram suas tarifas e taxas sobre grãos, carvão, aço e outros materiais em até 30\% entre 1880 e 1913 (Graff et al. 2014, 75). Economistas liberais, como F.Y. Edgeworth (1845–1926), tentaram “reafirmar”, como ele alegou, a teoria clássica da tributação como foi deixada por Mill. Usando mais uma vez seu famoso modelo de troca de duas pessoas, Edgeworth tentou determinar a incidência de impostos cobrados sobre mercadorias importadas em uma troca entre duas nações para concluir que “o imposto inflige mais perda em qualquer uma das partes, quanto menor for a elasticidade da demanda ou oferta dessa parte” (Edgeworth 1897, 48).\footnote{Como Edgeworth explicou, quando há circunstâncias monopolistas, o preço do bem tributado na verdade cairá. Esse resultado foi mais tarde chamado de “o paradoxo de Edgeworth”.}

Em suma, os Marginalistas e os primeiros economistas neoclássicos eram favoráveis à intervenção estatal como complementar às forças do mercado. Para aceitar “a violação do princípio geral de que o indivíduo sabe o que é melhor para si”, como afirmou Jevons (1882), não deve haver outra maneira de elevar a utilidade coletiva. Em questões sociais, os economistas neoclássicos, sendo observadores atentos de seu ambiente social, estavam particularmente atentos à melhoria das condições das classes trabalhadoras e à erradicação do pauperismo. Jevons reconheceu o problema como uma falha institucional e sugeriu mais educação como a única maneira de “elevar a inteligência e os hábitos previdentes do povo” (1883, 186; Sekerler e Sigot 2013). Da mesma forma, ele sugeriu bibliotecas públicas como um meio eficiente de aumentar a educação das pessoas (Mosselmans 2007, 81). Marshall também sublinhou a necessidade de uma melhor educação dos pobres (1890, 562. Cf. Jensen 1987). Por sua vez, Pigou (1920, 29–30) combinou a falta de educação dos pobres com sua baixa capacidade de prever o futuro e, portanto, com suas baixas economias, negligenciando assim o nível infeliz de seus salários. Tratamento análogo do pauperismo e dos baixos salários não é encontrado em Walker (1876) e Clark (1886), que acreditavam que os trabalhadores eram remunerados de forma justa de acordo com a produtividade marginal de seu trabalho.

\section*{Conclusão}

Em questões de política econômica, marginalistas e neoclássicos iniciais demonstraram uma continuidade notável com os economistas clássicos e mostraram heterogeneidade semelhante em suas recomendações políticas. Para começar com esta última, entre o liberalismo puro de Ricardo e Say e o radicalismo filosófico de Mill, havia conservadores como Malthus e Whigs como Senior, que mais ou menos concordavam com os princípios fundamentais da economia política e, no entanto, endossavam agendas políticas distintas. Análoga foi a diferenciação política de seus sucessores. De um lado, havia as posições ultraliberais de Pareto em linha com Clark e Fisher. Jevons, Edgeworth e mais tarde Pigou foram menos negativos quanto à intervenção estatal no caso de o mercado falhar em ser eficiente, marcando um liberalismo mais brando. Em contraste, Marshall e Walras, os dois maiores economistas desta era, embora adotando modelos teóricos separados, compartilhavam em comum muitas ideias dos planos utópicos socialistas de sua época. Propostas como a regulação dos lucros (Marshall) ou a nacionalização da terra (Walras) simplesmente enfureceriam Pareto ou Fisher, apesar de seus muitos parentescos teóricos. Assim, a primeira conclusão de nossa visão concisa das sugestões de política econômica da corrente principal até a década de 1920 é que elas eram muito heterogêneas, apesar de sua filiação inquestionável à escola do pensamento econômico Marginalista e Neoclássico Inicial.

Quanto à primeira característica, a continuidade nas agendas de política entre os Clássicos e os primeiros Neoclássicos, independentemente do deslocamento do escopo da análise econômica para questões microeconômicas de troca e produção, foi de fato um fenômeno transmissivo. Em questões de concorrência (tanto interna quanto externa), a mudança mais significativa veio do surgimento de grandes empresas que levaram alguns economistas como Marshall e Clark a apoiar políticas antitruste e a regulação estatal. A política monetária permaneceu quantitativista e mudou para o monetarismo apenas com Fisher. Finalmente, o papel do Estado na política pública interna continuou a ser uma questão disputável, da mesma maneira que na época anterior.

\section*{Bibliografia}

\begin{description}
    \item Backhouse, R.E. (2006), “Sidgwick, Marshall and the Cambridge School of Economics”, \textit{History of Political Economy}, 38 (1): 15–44.
    \item Backhouse, R.E., Baujard A. and Nishizawa, T. (2021), \textit{Welfare Theory, Public Action and Ethical Values}, Cambridge: Cambridge University Press.
    \item Bagehot, W. (1873), “Lombard Street: A Description of the Money Market,” in N. St John-Stevas (ed.) \textit{The Collected Works of Walter Bagehot}, London: The Economist, 1978 vol. IX, pp. 45–233.
    \item Bairoch, P. (1997), \textit{Victoires et déboires. Histoire économique et sociale du monde du XVI s. à nos jours}, Paris: Gallimard, 3 vols.
    \item Blaug, M. (1985), \textit{Economic Theory in Retrospect}, Cambridge: Cambridge University Press, 4th ed.
    \item Bruni, L. (2002), \textit{Vilfredo Pareto and the Birth of Modern Microeconomics}, Cheltenham: E. Elgar.
    \item Caldari, K. and Masini, F. (2011), “Pigouvian versus Marshallian Tax: Market Failure, Public Intervention and the Problem of Externalities”, \textit{European Journal of the History of Economic Thought}, 18 (5): 715–732.
    \item Clark, J.B. (1886), \textit{The Philosophy of Wealth: Economic Principles Newly Formulated}, New York: Macmillan.
    \item Clark, J.B. and Clark, J.M. (1901), \textit{The Control of Trusts}, Charleston, SC: Nabu Press 2010.
    \item Creedy, J. (1986), \textit{Edgeworth and the Development of Neoclassical Economics}, Oxford: Blackwell.
    \item Creedy, J. (1990), “Marshall and Edgeworth”, \textit{Scottish Journal of Political Economy}, 37 (1): 18–39.
    \item Dockès, P. (1996), \textit{La société n’est pas un pique-nique. Léon Walras et l’économie sociale}, Paris: Economica.
    \item Drakopoulos, S. (1989), “The Historical Perspective of the Problem of Interpersonal Comparisons of Utility”, \textit{Journal of Economic Studies}, 16 (5): 35–51.
    \item Edgeworth, F.Y. (1894), “Theory of International Values”, \textit{The Economic Journal}, 4 (16): 606–638.
    \item Edgeworth, F.Y. (1897), “The Pure Theory of Taxation”, \textit{The Economic Journal}, 7 (25): 46–70.
    \item Fawcett, H. (1872), \textit{Free Trade and Protection}, London: MacMillan, 6th ed. 1885.
    \item Fisher, I. (1911), \textit{The Purchasing Power of Money} (with H. Brown), New York: MacMillan, 1926.
    \item George, H. (1879), \textit{Progress and Poverty}, Dodo Press 2009.
    \item George, H. (1886), \textit{Protection or Free Trade}, London: Kegan Paul, Tresch and Co.
    \item Graff, M., Kenwood, A.G. and Lougheed, A.L. (2014), \textit{Growth of the International Economy, 1820–2015}, London: Routledge, 5th ed.
    \item Groenewegen, P. (1995), \textit{A Soaring Eagle: Alfred Marshall, 1842–1924}, Aldershot: E. Elgar.
    \item Groenewegen, P. (2003), “English Marginalism: Jevons, Marshall, and Pigou”, in W.J. Samuels, J.E. Biddle and J.B. Davis (eds.) \textit{A Companion to the History of Economic Thought}, Oxford: Blackwell, pp. 246–261.
    \item Heilbroner, R. (1953), \textit{The Worldly Philosophers}, New York: Simon and Schuster, 6th ed. 1992.
    \item Hobsbawm, E. (1987), \textit{The Age of Empire, 1875–1914}, London: Weidenfeld and Nicolson.
    \item Jensen, H.E. (1987), “Alfred Marshall as a Social Economist”, \textit{Review of Social Economy}, 45 (1): 14–36.
    \item Jevons, W.S. (1865), \textit{The Coal Question}, New York: Kelley 1965.
    \item Jevons, W.S. (1871), \textit{The Theory of Political Economy}, London: MacMillan 1879, 2nd ed.
    \item Jevons, W.S. (1882), \textit{The State in Relation to Labour}, London: MacMillan 1968.
    \item Jevons, W.S. (1883), \textit{Methods of Social Reform}, New York: Kelley reprints 1965.
    \item Keynes, J. M. (1924), “Alfred Marshall 1842–1924”, \textit{Economic Journal}, 34 (135): 311–372.
    \item Laidler, D. (1991), \textit{The Golden Age of the Quantity Theory. The Development of Neoclassical Monetary Economics, 1870–1914}, New York: Allan.
    \item Legris, A. (2000), “La distribution des revenus chez Walras et Pareto: une analyse Comparative”, in P. Dockès et al. (eds.) \textit{Les Traditions économiques françaises}, Paris: CNRS éditions, pp. 503–518.
    \item Lewis, A.W. (1978), \textit{Growth and Fluctuations 1870–1913}, London: Allen \& Unwin.
    \item Marshall, A. (1890), \textit{Principles of Economics}, London: MacMillan Press 1985, 8th ed. 1920.
    \item Marshall, A. (1919), \textit{Industry and Trade}, London: MacMillan Press, 3d ed. 1920.
    \item Marshall, A. (1923), \textit{Money, Credit and Commerce}, London: MacMillan Press.
    \item Medema, S. (2009), \textit{The Hesitant Hand: Taming Self-Interest in the History of Economic Ideas}, Princeton, NJ: Princeton University Press.
    \item Mill, J.S. (1848), \textit{Principles of Political Economy with Some of Their Applications to Social Philosophy}, New York: A. Kelley, 1973.
    \item Mosselmans, B. (2007), \textit{William Stanley Jevons and the Cutting Edge of Economics}, London: Routledge.
    \item Pareto, V. (1897), “The New Theories of Economics”, \textit{Journal of Political Economy}, 5: 485–502.
    \item Pareto, V. (1900), \textit{Libre-échangisme, protectionnisme et socialisme}, in G. Busino (ed.) \textit{Œuvres Complètes} vol. 4 Genève: Librairie Droz 1965.
    \item Pareto, V. (1909), \textit{Manuel d’Economie Politique}, in G. Busino (ed.) \textit{Œuvres Complètes} vol.7, Genève: Librairie Droz 1966.
    \item Pareto, V. (1916), \textit{Traité de Sociologie générale}, in G. Busino (ed.) \textit{Œuvres Complètes} vol. 12, Genève: Librairie Droz 1968.
    \item Peart, S. (1996), \textit{The Economics of William Stanley Jevons}, London: Routledge.
    \item Pigou, A.C. (1906), \textit{Protective \& Preferential Import Duties}, London: MacMillan.
    \item Pigou, A.C. (1920), \textit{The Economics of Welfare}, London: MacMillan, 3d ed. 1929.
    \item Pigou, A.C. (1921), \textit{The Political Economy of War}, London: MacMillan.
    \item Pigou, A.C. (1925), \textit{Memorials of Alfred Marshall}, London: MacMillan.
    \item Piketty, T. (2013), \textit{Le Capital au XIX siècle}, Paris: Seul.
    \item Potier, J.-P. (2011), “Marché du travail et législation sociale dans la pensée de Léon Walras”, \textit{Œconomia}, 1–3: 437–458.
    \item Potier, J.-P. (2019), \textit{Léon Walras, économiste et socialiste libéral}, Paris: Classiques Garnier.
    \item Pressman, S. (2014), \textit{Fifty Major Economists}, London: Routledge, 3d ed.
    \item Raffaelli, T. (2003), \textit{Marshall’s Evolutionary Economics}, London: Routledge.
    \item Raffaelli, T., Becattini, G. and Dardi, M. (2006), \textit{The Elgar Companion to Alfred Marshall}, London: Routledge.
    \item Robbins, L. (1973), \textit{The Evolution of Modern Economic Theory}, London: MacMillan Press.
    \item Schumpeter, J.A. (1954), \textit{History of Economic Analysis}, London: Allen \& Unwin.
    \item Sekerler, P. and Sigot, N. (2013), “William Stanley Jevons et la « réforme sociale »: Une théorie du bien-être sans postérité”, \textit{Cahiers d’économie Politique}, 64: 221–251.
    \item Sidgwick, H. (1883), \textit{The Principles of Political Economy}, London: MacMillan and Co., 3d. ed. 1901.
    \item Steiner, P. (1995), “Vilfredo Pareto et le protectionnisme: l’économie politique appliquée, la sociologie générale et quelques paradoxes”, \textit{Revue Economique}, 46 (5): 1241–1262.
    \item Tobin, J. (1985), “Neoclassical Theory in America: J.B. Clark and Fisher”, \textit{American Economic Review}, 75 (6): 28–38.
    \item Walker, F.A. (1876), \textit{The Wages Question. A Treatise on Wages and the Wages Class}, New York: H. Holt \& Co.
    \item Walker, F.A. (1884), \textit{Political Economy}, New York: H. Holt \& Co.
    \item Walker, D.A. (2003), “Early General Equilibrium Economics: Walras, Pareto, and Cassel”, in W.J. Samuels, J.E. Biddle and J.B. Davis (eds.) \textit{A Companion to the History of Economic Thought}, Oxford: Blackwell, pp. 278–293.
    \item Walras, L. (1874/1877), \textit{Eléments d’économie politique pure ou Théorie de la richesse sociale}, in Auguste et Léon Walras \textit{Œuvres économiques complètes}, vol. VIII, Paris: Economica 1988.
    \item Walras, L. (1896), \textit{Etudes d’économie sociale: Théorie de la répartition de la richesse sociale}, in Auguste et Léon Walras \textit{Œuvres économiques complètes}, vol. IX, Paris: Economica 1990.
    \item Walras, L. (1898), \textit{Etudes d’économie appliquée: Théorie de la production de la richesse sociale}, in Auguste et Léon Walras \textit{Œuvres économiques complètes}, vol. X, Paris: Economica 1992.
    \item Wicksell, K. (1898), \textit{Interest and Prices: A study of the Causes Regulating the Value of Money}, London: Macmillan 1936.
    \item Wicksell, K. (1906), \textit{Lectures on Political Economy}, 2 vols, London: Routledge 1938.
    \item Wicksell, K. (1907), “The Influence of the Rate of Interest on Prices”, \textit{The Economic Journal} 17 (6): 213–220.
    \item Wicksteed, P.H. (1910), \textit{The Common Sense of Political Economy, Including a Study of the Human Basis of Economic Law}, London: MacMillan.
    \item Winch, D. (1969), \textit{Economics and Policy: A Historical Study}, London: Hodder and Stoughton.
    \item Zouboulakis, M. (2003), “L’équation de l’échange de Fisher et la loi des gaz parfaits. Les limites de l’analogie physique en économie”, \textit{Economies et Sociétés}, Série PE, 33: 2191–2206.
    \item Zouboulakis, M.S. (2014), \textit{The Varieties of Economic Rationality: From Adam Smith to contemporary Behavioral and Evolutionary Economics}, London: Routledge.
\end{description}

\end{document}